% !TeX root = ../papers/nc-p3.tex
\section{種数二の滑らかな五次点グラフの偏極}
\label{sec:genus-two-quintic-ruling}
\begin{theorem}\label{thm:genus-two-quintic-ruling}
四辺の左右射影ファイバーが空か一点となる標準有限Frobenius候補で、
点グラフの一次元部分が滑らかな整五次曲線 $C$ で、孤立残余と交わらないとする。
種数は1または2である。種数2では孤立残余の長さは2で、標準四移動は超楕円対合となる。
$H=\omega_C$ とし、添字を固定すると、偏極は二つの相異なるWeierstrass点 $a,b$ に対する
$H^2(a),H^2(b)$ を交互に繰り返す。各埋込みを含む二次曲面は錐である。
\end{theorem}
\begin{proof}
平面五次なら種数6であり全位置で平面となるが、前節の平面積の完全交叉次数3に反する。
非退化ならCastelnuovo上界が2、系\ref{cor:point-graph-excess}の孤立残余長は $2g-2$ である。
以下 $g=2$ とする。次数5の $L_i$ は完全な四切断を持ち、$h^0(L_i^2)=9$ だから曲線は二次曲面 $Q_i$ に入る。
非退化整曲線の次数が4を超えるのでこの曲面は一意かつ既約である。
$L_i=L_{i+1}$ なら全関係がその二次曲面の対角で零となるため、隣接偏極は相異なる。

$Q_2$ が滑らかなら曲線の型は $(3,2)$ または逆で、二つの射影の次数は2と3である。
錐ならその解消 $\mathbb F_2$ の負切断を $s$、ファイバーを $f$ とする。
曲線の固有変換は元の滑らかな曲線と同型である（頂点イデアルの曲線への制限は可逆）。
その類は $ms+5f$、種数は $-m^2+5m-4$ で $5-2m\ge0$ なので $m=2$ となる。
$s$ は一点 $t$ で切り、超平面 $s+2f$ の制限は $L_2=H^2(t)$ である。
どちらの曲面でも次数2の鉛筆は種数2曲線唯一の超楕円鉛筆である。

一般の $p\in C$ を固定する。反対二点縮約を三行三列に書いた行列の三次行列式 $S_p$ は $C$ で零になる。
その $Q_2$ への制限が非零であることを示す。恒等零なら曲面の一般点で行列の階数は2である
（曲線の一般点ですでに2）。二次小行列からなる随伴行列の各成分を
滑らかな二次曲面または錐の解消のCox環に書く。この環は多項式環である。
階数1ゆえ、共通因子を除いた列と行の外積と一つの同次共通因子に分解できる。
実際、一列の最大公約元を除き、他の列をその有理函数倍に書くと分母は単元となる。
得た列と行は各々同じ類の三切断からなり、それぞれ固定因子を持たない。

これら二写像は $C$ 上で、$j_1(C),j_3(C)$ をそれぞれ $j_1(p),j_3(\sigma^{-1}p)$ から射影したものを与える。
核直線は一般点で一意であり、これは定理\ref{thm:opposite-curve-divisor}の左右核の記述による。
両射影は平面を張り、一般の $p$ では像は円錐曲線でもない。実際、射影束は次数4で、
円錐曲線なら唯一の二重被覆から $H^2$ を引き戻した系、すなわち $L_1(-p)=H^2$
（他方も同様）である。これは高々一点に限られる。

滑らかな二次曲面では、片方の次数が零の切断列は一つの射影を経由する。
その次数は小行列の二次数 $(2,2)$ により高々2で、像が平面を張れば円錐曲線になる。
従って列と行は共に両次数が正で、双方の型は $(1,1)$ である。
錐の解消でも、固定因子のない系の類 $u s+v f$ は $v\ge2u$ を満たす。
$u=0$ の場合は射が超楕円射影を経由し、その像の次数は元の射影束の次数4から高々2である。
両方で $u\ge1$ なので、小行列の類 $2s+4f$ から双方とも $s+2f$ になる
（共通因子の有効類の二係数も非負）。
従っていずれの曲面でも各射影は三つの超平面切断で与えられる。
特に曲線への制限で次数5から4へ下がり、ある点 $x$ で $L_2(-x)=L_1(-p)$ である。
これは一般の $p$ に成り立たない。実際 $L_2/L_1$ は非自明で、この束を $x-p$ と書く方法は有限である。
一つの表示 $x_0-p_0$ と比較すれば $x+p_0\sim x_0+p$、
二因子が異なるならどちらも超楕円ファイバーである。この矛盾で $S_p|_{Q_2}\ne0$ を得る。

曲面上でこの三次切断から $C$ を除くと、滑らかな場合は次数3で $C$ を切る一本の母線、
錐の解消では $s+f$ の有効因子が残る。後者は $s$ と一本のファイバーの和である。
どちらの場合も残余の $C$ への制限 $E_p$ は次数3で、束は $P_2=L_2/H$ となる。
$D_p$ を前出の定理の長さ2の階数低下因子とする。$S_p$ の微分は二次小行列の線型結合なので
$D_p$ で零である。一般の $p$ では $D_p$ は錐の頂点を避ける（両射影は有限）。
滑らかな曲面で局所的に $S_p=c e$、$c=0$ が滑らかな $C$ を切るため、
制限微分は $e\,dc$ であり $D_p\subset E_p$ がスキームとして従う。
従って残余は一点 $r(p)$ であり、束として
\[
 \mathcal O(r(p))=P_2\otimes K\otimes\mathcal O(-p-\tau p),\qquad \tau=\sigma^{-1}.
\]
右辺は全ての $p$ で次数1のAbel像に入る（その像は閉）。ゆえに $r:C\to C$ は射である。

$r$ が非定数ならHurwitzにより自己同型で、Jacobian上で $r_*=-1-\tau_*$ となる。
不変Hermite内積で両者がユニタリなので $\tau_*+\tau_*^{-1}=-1$、従って $r_*=\tau_*^2$ である
（有限個の係数を複素数に埋めればよい）。自己同型のJacobian作用は忠実なので
$\tau$ は位数3で $r=\tau^2$ となる。どの位置でも同じJacobian恒等式を適用でき、
右辺の準同型が非零なので残余写像は全て同じ $\tau^2$ になる。
上式を各位置で比較すると全ての $P_i$ が一致する。これは隣接偏極が異なることに反する。

ゆえに $r$ は定数で $\tau_*=-1$、$\tau$ は超楕円対合である。
系\ref{cor:curve-divisor-monodromy}により $K=1$。上の式より $P_i=H(t_i)$、
$L_i=H^2(t_i)$ が全位置で従う。滑らかな二次曲面の場合に $P_i$ は基点自由な第三次鉛筆だが、
$H(t_i)$ は $t_i$ を基点に持つので全位置が錐である。
一階差 $d=L_{i+1}/L_i$ は $K=1$ により一定、$\sigma^*d=d$ かつ $\sigma^*d=d^{-1}$ なので $d^2=1$。
従って $L_{i+2}=L_i$ で、$L_{i+4}=\sigma^*L_i=L_i$ より各 $t_i$ はWeierstrass点となる。
\end{proof}

\begin{theorem}\label{thm:genus-two-quintic-exclusion}
前定理の種数2の場合は存在しない。従って滑らかな整五次曲線と交わらない有限残余からなる
点グラフの候補は、楕円五次曲線そのものに限られる。
\end{theorem}
\begin{proof}
超楕円座標 $x$ の零と極を $a,b$ の下に取ると
$y^2=x F(x)$、$F$ は四次で零点は互いに異なり $F(0)\ne0$ となる。
$H=\mathcal O(2b)$ とすると $L_0=L_2=H^2(a)$ の基底は $1,x,x^2,y/x$、
$L_1=L_3=H^2(b)$ の基底は $1,x,x^2,y$ である。これらを添字 $0,1,2,3$ で表す。
第四の文字を奇文字と呼ぶ。偶数位置で $a$ は第四座標点、$b$ は第三座標点である
（$b$ で $x^2$ の極は位数4、$y/x$ の極は位数3）。
標準四移動は $y\mapsto-y$ なので、四番目の辺での位置0の座標同一視は
奇文字だけの符号変更である。偏極は完全で曲線は非退化なので射影同一視はこれで定まる。

四次Frobeniusテンソルを $w$ とする。他の二添字を固定した係数は
隣接関係に属するので、隣接二切断の積で評価して零になる。巡回する最後の辺では上の符号変更を施す。
$y$ を一個含む関数だけを取れば、偶数から奇数位置への辺で
\[
 w_{30}=w_{23}=0,\qquad w_{3,j+1}=-w_{j,3}\quad(j=0,1)
\]
が成り立つ（他の位置を略記する）。逆の型の辺では転置した式を使う。
位置3から0の辺でのみ交換式の負号が正号になる。これらは $y$ の奇部分での
$x^{-1},1,x,x^2$ の独立性だけを用い、$F$ の係数を選ばない。

奇文字が一個か三個ならその四添字の係数は零である。
一個の場合は奇文字を右に一つずつ、他の三文字と12回交換する。
交換不能なら端の零式により係数が零である。全交換が可能なら、他の各文字との交換は
3ステップおきに4回起き、数値添字の増減は辺の偶奇に従い打ち消し合う。
三個の場合は一つだけの通常文字を右へ12回交換し、毎回交互にその値が増減する。
どちらも全添字が元へ戻り、最後の辺を3回、他の辺を計9回通るので符号は負となる。
従って標数零では消滅する。
特に $w_{3,j,2,k}=0$ が全 $j,k$ で成り立つ。$j,k$ が通常文字なら奇文字は一個、
$j=3$ なら位置1から2で転置した端の零式、$k=3$ なら位置2から3で $w_{23}=0$ を使う。
これらの係数消滅は\path{computations/verify_genus_two_quintic_ruling.py}でも有理数上の行空間で検算する。

よって $j_0(a)$ での三次縮約は、中位置で $j_2(b)$ が消す三平面 $U_2$ に支持される。
左端は $j_1(a)$ が消す三平面 $U_1$ に支持され、右端支持は次元3である。
実際、隣接切断の階数6と左右の点射影ファイバーの一意性により、隣接関係の評価の像は三平面全体である。
従って末尾の縮約は $R_1$ に入る $U_1\otimes U_2$ の三次元部分空間を与える。
$U_1,U_2$ はそれぞれ $L_1(-a),L_2(-b)$ の全切断である。
$2a\sim2b$ より両束は共通の次数4の束 $N$ となる。この平面写像の像は平面を張り、二次曲線に入らない。
そうなら唯一の二重被覆から $N=H^2$ となるが $a\ne b$ に反する。
従って $U_1\otimes U_2$ の全双一次関係はこの同一視のもとで三つの交換子であり、$R_1$ は全てを含む。

$N(a+b)=L_1(b)=L_2(a)$ は次数6の非常に豊富な束で、五つの全切断を持つ。
二つの四切断空間の共通部分は上の三切断であり、和は全切断となる。
対応する射影四空間でこれら三座標を $h$、二つの補座標を $u,v$ と書く。
三つの交換子以外の三関係は三つの二次式になる。
$h\ne0$ では点 $[h:u:v]$ と共通座標が比例かつ非零な二点の組とは一対一に正則に復元される。
したがってこの開集合上で三式の共通零点の一次元の台は、次数6の $C$ に限られる。
$h=0$ の割線では全二次式は $uv$ の倍数で、どれかの係数が非零である。
全て零なら $(j_1(a),j_2(b))$ が点対応に入り、$a\ne b$ かつ $j_1(a)$ の次点の一意性に反する。
従って割線も成分ではない。三式は余次元3の完全交叉を定め、次数8である。
その一次元成分の台は次数6の曲線一つなので次数は6の倍数となるはずで矛盾する。
\end{proof}
